% ============================================================================
% DRAFT — OCEANS 2026 companion paper (behaviour-profiling pipeline).
% Sibling paper: conference_101719.tex (protest-behavior -> lighting).
% TODO before submission:
%   - finalise author list / affiliations (placeholder below, mirrors the consortium)
%   - generate figures listed in comments (behaviour dashboard charts, seg before/after,
%     pipeline diagram); verify all \includegraphics paths
%   - verify every citation's details
%   - decide final length / trim to target track (full vs work-in-progress)
%   - build the human-verified test sets before making any "A beats B" claim (see Sec. Limitations)
% ============================================================================
\documentclass[conference]{IEEEtran}
\IEEEoverridecommandlockouts
\usepackage{cite}
\usepackage{amsmath,amssymb,amsfonts}
\usepackage{algorithmic}
\usepackage{graphicx}
\usepackage{textcomp}
\usepackage{xcolor}
\usepackage{enumitem}
\usepackage{booktabs}
\def\BibTeX{{\rm B\kern-.05em{\sc i\kern-.025em b}\kern-.08em
    T\kern-.1667em\lower.7ex\hbox{E}\kern-.125emX}}
\begin{document}

\title{From Footage to Ethogram: An Automated Vision--Language Pipeline for
Quantifying Behaviour and Affective State of a Captive Octopus}

% TODO: confirm/finalise author list and ordering with the consortium.
\author{
\IEEEauthorblockN{
Siddharth Raj\IEEEauthorrefmark{1},
Harald Burgsteiner\IEEEauthorrefmark{2},
Wolfgang Slany\IEEEauthorrefmark{3}
}
\IEEEauthorblockA{\IEEEauthorrefmark{1}
Research and Development, Arishna IoT Solutions, Bangalore, India}
\IEEEauthorblockA{\IEEEauthorrefmark{2}
Institute for Digital Media Education, University of Teacher Education Styria, Graz, Austria}
\IEEEauthorblockA{\IEEEauthorrefmark{3}
Institute of Software Engineering and Artificial Intelligence, Graz University of Technology, Graz, Austria}
}
\maketitle

\begin{abstract}
Continuous camera monitoring of captive marine animals produces far more video
than caretakers can review, yet most of it is used only reactively. We present
an automated vision--language pipeline that converts months of continuous
multi-camera aquarium footage of a single \emph{Octopus vulgaris} into a
quantified behavioural and affective-state time-series. The pipeline detects
when the animal is visible (a frozen CLIP encoder with a lightweight probe under
aspect-preserving preprocessing), extracts short behaviour clips gated on
visibility and motion, and re-prompts a large vision--language model to emit a
structured behavioural record per clip (behaviour class, posture, activity,
location, context and body colour). Aggregating $3{,}083$ analysed clips yields
an interpretable activity budget, an exposure-normalised circadian profile with a
late-afternoon activity peak, and a clear stimulus response in which human
presence nearly doubles movement and raises a transparent arousal index. To make
the system deployable and self-hosting, we distil the large teacher into small
local models: a captioning student (Qwen3-VL-2B, QLoRA) that runs 4-bit on a
laptop, and a single-class octopus mask model auto-labelled by GroundingDINO and
SAM2. The segmentation presence signal reliably rejects tank-glass reflections,
a dominant false-positive source. We frame results in terms of arousal and
behavioural state rather than emotion, report the pipeline's current
limitations honestly, and position the work as the profiling layer complementary
to real-time single-behaviour response systems.
\end{abstract}

\begin{IEEEkeywords}
Computer Vision; Vision--Language Models; Behaviour Analysis; Marine Animal
Welfare; Model Distillation; Edge AI; Cephalopods
\end{IEEEkeywords}

\section{Introduction}
Aquarium and controlled-habitat monitoring increasingly relies on
continuous camera coverage, but the resulting video vastly exceeds what human
caretakers can review. Most footage is therefore consulted only reactively,
after an event, and the rich record of an animal's day-to-day behaviour is left
unquantified. For welfare-oriented husbandry the more valuable question is
longitudinal: \emph{how does this animal spend its time, how does its rhythm
look, and how does it respond to its environment?}

Octopuses make this both important and hard. They are highly intelligent,
behaviourally complex invertebrates whose welfare is a growing concern, yet their
non-rigid bodies, rapid posture changes, active camouflage, frequent denning, and
the reflections and floating debris of tank environments defeat naive visual
monitoring~\cite{disc_o}. A companion line of work from our consortium targets a
single, actionable behaviour---a ``protest'' limb gesture used to request tank
lighting---and closes the loop by automatically actuating the lights. That system
is reactive and behaviour-specific by design.

This paper addresses the complementary problem: \emph{understanding the whole
behavioural repertoire over time}. We present an automated pipeline that turns
continuous multi-camera footage of one captive \emph{O.\ vulgaris} (``Nity'')
into a structured behavioural and affective-state time-series, together with the
small local models needed to run it cheaply and offline. Our contributions are:
\begin{enumerate}[leftmargin=*, itemsep=1pt, topsep=2pt]
\item An end-to-end pipeline from raw footage to a per-clip \emph{structured
behavioural record}, using a large vision--language model (VLM) as a
zero-shot behavioural annotator with snap-to-list validation.
\item An aggregate behavioural analysis of the animal over several months: an
interpretable activity budget, an exposure-normalised circadian profile, and a
quantified stimulus response, framed as arousal / behavioural state.
\item Teacher$\rightarrow$student distillation into deployable local models---a
4-bit captioning student that runs on a laptop, and a single-class octopus
segmentation model auto-labelled by GroundingDINO$+$SAM2 whose mask-area signal
robustly rejects reflections.
\item An honest account of current limitations---most importantly the absence of
human-verified held-out test sets---and the diverse-footage data strategy that
addresses the dominant generalisation gap.
\end{enumerate}

\section{Related Work}
\textbf{Cephalopod behaviour and welfare.} Interest in octopus cognition and
welfare has grown alongside their inclusion in animal-welfare frameworks. Prior
consortium work introduced a framework for decoding inter-species communication
with octopodes and a real-time detector for a light-request behaviour that
actuates the habitat~\cite{disc_o}. Our work is the profiling counterpart to that
response system.

\textbf{Foundation models as annotators.} General-purpose VLMs can describe and
classify images zero-shot, and open-vocabulary detectors and promptable
segmenters (GroundingDINO~\cite{groundingdino}, SAM/SAM2~\cite{sam2}) enable
labelling without task-specific supervision. We use these as \emph{teachers}: an
expensive model produces structured labels that we distil into small
students~\cite{lora}. CLIP~\cite{clip} provides the frozen visual features for our
lightweight presence detector; the deployed segmenter uses an efficient
LR-ASPP/MobileNetV3 head~\cite{mobilenetv3}.

\section{System and Methods}
The tank is observed by a multi-camera rig (6 IP cameras, 4K, infrared and colour
night vision) with a manually selected region of interest around the den for each
camera; the same installation supports the companion response system. The
pipeline (Fig.~\ref{fig:pipeline}) has four stages.

% FIGURE: pipeline diagram (detection -> extraction -> structured extraction -> distillation).
% generate from src/ (draw in draw.io / tikz). Placeholder path below.
\begin{figure}[t]
    \centering
    \includegraphics[width=0.48\textwidth]{assets/pipeline_behaviour.pdf}
    \caption{Footage-to-ethogram pipeline: visibility detection $\rightarrow$
    clip extraction $\rightarrow$ VLM structured behavioural extraction
    $\rightarrow$ distillation into local caption and segmentation students.}
    \label{fig:pipeline}
\end{figure}

\subsection{Visibility Detection}
A frozen CLIP ViT-B/32 encoder feeds a small multi-layer-perceptron probe
(512$\rightarrow$256$\rightarrow$64$\rightarrow$2) trained to classify
\emph{octopus-visible} vs.\ \emph{hidden}. The single most important design choice
was preprocessing: CLIP's default centre-crop discards $33$--$44\%$ of a 16:9
frame, and \emph{aspect-ratio mismatch}---not backbone capacity---was the root
cause of poor field performance. Replacing the crop with aspect-preserving
letterbox padding, and mining $\sim$1.7k reflection/infrared-noise hard
negatives, gave a probe that reaches $96.8\%$ on our internal detection test set
while sharply reducing confident false positives on reflections.

\subsection{Clip Extraction}
For each source video we compute, per second, the probe's visibility probability
and an absolute changed-pixel motion fraction (with the burned-in timestamp region
masked so the ticking clock is not counted as motion). A non-overlapping
$20\,\mathrm{s}$ window is kept when the octopus is visible for a majority of the
window and mean motion exceeds a threshold. We deliberately use the
\emph{absolute} changed-pixel fraction: per-video normalised motion incorrectly
promotes static scenes (e.g.\ a flickering lamp) to ``motion''.

\subsection{Structured Behavioural Extraction}
Rather than treating a caption as the end product, we treat it as one field of a
\emph{structured behavioural record}. For each extracted clip we select the
clearest frames (contrast-enhanced for the dim infrared footage, scored by the
visibility probe) and prompt a large VLM (Qwen3-VL-235B, accessed via API) to
return a JSON record: \{present, behaviour (7-class ethogram), posture, activity
level, location, context, body colour, colour/texture change, confidence\}. Each
field is validated by snapping to an allowed value list; colour fields are gated
per clip, so infrared/greyscale clips (detected by low BGR channel divergence) are
instructed to leave colour \emph{uncertain}. The run over $3{,}205$ clips cost
$\approx\$2.22$ ($\sim\$0.0006$/clip). We note---and exploit---that the VLM is a
far stronger presence filter than the detector: on one batch it judged $63\%$ of
detector-``verified'' clips to actually contain no octopus, almost all tank-glass
reflections.

\subsection{Distillation to Local Models}
\label{sec:distill}
A central design principle is that the expensive foundation models are used only
as \emph{offline teachers}; nothing large is deployed. Their outputs become
training data for small students, so the running system is self-hosting and needs
no GPU server or paid API. We apply the same teacher$\rightarrow$student recipe
twice: (i) the VLM's captions distil into a compact captioning model, and (ii) the
GroundingDINO$+$SAM2 masks distil into a single-class segmentation model
(Sec.~\ref{sec:models}). In both cases the teacher provides the labels a small
model could never be trained on directly, and the student is deliberately kept as
small as the task allows---a single class, low input resolution---so that
continuous in-situ inference is feasible on commodity hardware.

\subsection{Diverse-Footage Harvesting}
\label{sec:harvest}
Early models were trained on footage from a single week, which we found to be the
binding generalisation limit. We therefore built a streaming harvester that crawls
the footage archive, and for each candidate video \emph{probes} a handful of
frames before committing to a full scan---skipping empty ``den'' videos cheaply.
Because the archive server imposes a low aggregate bandwidth cap (measured
$\sim$5\,MB/s, with parallel connections yielding only $1.6\times$), the harvester
runs as a single CPU worker that streams rather than bulk-downloads, keeping only
short clips and a detailed coverage ledger so that skipped footage can be revisited
later. This expanded coverage from $7$ days to $>\!100$ distinct days.

\section{Behavioural and Affective-State Analysis}
Aggregating the $3{,}083$ present clips (joined to absolute clock time) yields an
interpretable behavioural profile. We stress a framing choice: we report
\textbf{arousal and behavioural state}, not emotion. Our transparent arousal index
is $0.6\cdot\mathrm{activity} + 0.4\cdot\mathrm{posture\text{-}spread}$, with
static body colour as a secondary channel available only on colour cameras.

% FIGURE: activity budget + circadian + stimulus-response (from data/behaviour_dashboard.html)
\begin{figure}[t]
    \centering
    \includegraphics[width=0.48\textwidth]{assets/behaviour_findings.pdf}
    \caption{Aggregate behaviour: (a) activity budget; (b) exposure-normalised
    circadian visible-activity rate; (c) motion and arousal with vs.\ without human
    presence. Generated from the structured records.}
    \label{fig:findings}
\end{figure}

\textbf{Activity budget.} Exploration/manipulation $41\%$, resting $33\%$,
human/enrichment interaction $14\%$, reaching out of water $9\%$, crawling $2\%$,
swimming/jetting $1\%$.

\textbf{Circadian structure.} Normalising present-clip counts by all extracted
windows per hour (an exposure correction) reveals a strong diel pattern: visible
activity is $\sim$1--5\% overnight, rises to a $\sim$45\% peak at 17:00 with a
13:00--19:00 plateau, plus a dawn bump around 05:00--06:00.

\textbf{Stimulus response.} Comparing clips with vs.\ without a human present, mean
motion nearly doubles ($0.045\rightarrow0.095$) and the arousal index rises
($0.46\rightarrow0.68$). On colour cameras, the dark red-brown body colour is most
common at baseline ($\sim$16\%) and less so during human interaction ($\sim$6\%),
consistent with a colour/posture shift under stimulus.

We regard the \emph{contrasts} (e.g.\ human-present vs.\ absent) as the robust
findings; absolute levels depend on the upstream presence gate and will shift as
the detector improves.

\section{Local Models}
\label{sec:models}
\subsection{Captioning Student}
We distil the teacher's one-sentence descriptions into Qwen3-VL-2B using QLoRA
(rank 16). On a $50$-clip held-out split, embedding similarity to the teacher
improves from $0.702$ (base model) to $0.834$ and ROUGE-L from $0.269$ to
$0.455$. Merged and quantised to 4-bit with MLX, the student is $\sim$1.7\,GB and
produces a caption in $\sim$3\,s on a 16\,GB laptop with no GPU; a cross-platform
Transformers backend covers non-Apple hardware.

\subsection{Segmentation Student}
Following the distillation principle of Sec.~\ref{sec:distill}, we target the
\emph{smallest} model that yields usable masks: a single foreground class at low
resolution, with GroundingDINO$+$SAM2 acting purely as an offline auto-labeller,
never deployed. To obtain masks without hand labelling, the teacher pairs
GroundingDINO (box) with SAM2 (mask): the most-confident frame seeds SAM2 video
propagation, and largest-blob plus area-continuity cleanup remove transient
errors. This corrected the two systematic failures of per-frame prompting
(infrared bright-tool bleed, colour background bleed) and, importantly, rejects the
reflection camera by low seed confidence. Auto-labelling produced $4{,}412$
(image, mask) pairs, from which we distilled candidate students spanning a
size sweep (a compact U-Net at several widths and an LR-ASPP head); we retain the
smallest model clearing our quality bar, an LR-ASPP student of $3.2$\,M
parameters.

Two results, one positive and one cautionary. \emph{Positive:} used as a
\emph{presence} signal, the mask area is highly reflection-robust. A model trained
with empty-mask negatives reaches AUC $0.86$ overall and $0.99$ against
reflections, with $0\%$ reflection false-positives at an area threshold while
retaining $88\%$ present-recall---directly attacking the pipeline's dominant
false-positive source. \emph{Cautionary:} pixel-accurate masks are not yet
solved---mask IoU plateaus near $0.47$ across model sizes, architectures and
augmentation. We diagnosed this as a \emph{video-diversity} generalisation gap
(only $62$ distinct training videos; training IoU $0.68$ vs.\ validation $0.47$,
failing by mislocating an otherwise correctly-sized blob), which motivates the
harvesting effort of Sec.~\ref{sec:harvest} rather than a larger model.

% FIGURE: seg before/after (data/segmentation_demo) + presence-gate ROC
\begin{figure}[t]
    \centering
    \includegraphics[width=0.48\textwidth]{assets/segmentation_demo.pdf}
    \caption{Auto-labelling teacher (GroundingDINO$+$SAM2) masks tracked across a
    clip on a colour camera (top) and an infrared camera (bottom). SAM2 video
    propagation with largest-blob and area-continuity cleanup yields the temporally
    consistent masks used to train the segmentation student; the reflection camera
    is rejected by low seed confidence.}
    \label{fig:seg}
\end{figure}

\section{Relevance to the OCEANS Community}
The work targets intelligent sensing and longitudinal behavioural analysis for
controlled marine habitats. It aligns with OCEANS themes in imaging and vision
systems and in welfare-oriented monitoring, and it emphasises \emph{edge
deployability}: the distilled students run offline on commodity hardware, avoiding
cloud dependence for continuous in-situ monitoring. Where the companion system
acts on a single behaviour in real time, this pipeline supplies the quantitative
behavioural baseline against which such events---and future welfare
interventions---can be interpreted.

\section{Limitations and Ongoing Work}
We are deliberate about what is not yet established.
\begin{itemize}[leftmargin=*, itemsep=1pt, topsep=2pt]
\item \textbf{No human-verified held-out test sets yet.} All reported numbers are
on internal evaluation splits (teacher-derived or auto-labelled). We therefore make
\emph{no} claim that one model beats another; building small human-verified test
sets for detection, captioning and segmentation is the immediate next step and a
prerequisite for any comparative claim.
\item \textbf{Presence-gate noise propagates.} Absolute behavioural rates depend on
the detector; we report contrasts, which are robust to a constant gate bias.
\item \textbf{Context label is coarse.} The ``enrichment interaction'' context
fires whenever a permanent tank object is visible ($\sim$66\% of clips); the clean
stimulus contrast is human-present vs.\ absent.
\item \textbf{Segmentation.} Pixel-accurate masks await more distinct training
videos; infrared masks additionally need a prompt-refinement fix. A static-feature
behaviour classifier failed (below the majority baseline), confirming that
behaviour classification needs temporal/motion features rather than more labels.
\end{itemize}
Ongoing work follows two tracks. First, we are lifting the segmentation student's
mask accuracy \emph{by data rather than by size}: a recently completed diverse
harvest (Sec.~\ref{sec:harvest}) expands training coverage from $7$ to over $100$
distinct days, and---consistent with the video-diversity diagnosis---we are
retraining the same small model on this broader set and re-selecting the smallest
architecture that meets the quality bar, so the deployed model stays compact while
its masks improve. Second, we distil the structured behavioural extractor itself
into the local $2$-B student (JSON targets in place of a sentence). Both tracks
depend on the human-verified test sets above, which we are building in parallel.

\section{Conclusion}
We presented an automated vision--language pipeline that converts continuous
octopus footage into a quantified behavioural and affective-state time-series, and
distils its large teachers into local models that run offline. The aggregate
analysis surfaces an interpretable activity budget, a clear circadian rhythm and a
measurable stimulus response, while the reflection-robust presence signal improves
the extraction itself. Framed as arousal and behavioural state and reported with
its limitations, the system offers a practical, welfare-oriented layer for
longitudinal behavioural understanding in controlled marine habitats.

\section*{Ethical Statement and Data Availability}
The work was conducted in a private aquarium, with only sub-threshold interactions
w.r.t.\ EC directive 2010/63/EU of several caretakers with a single
\emph{O.\ vulgaris} (supplier: Flying Sharks, Portugal). Husbandry and enrichment
practices followed best practices for octopus aquariums in compliance with FELASA
guidelines. The video data used in the study are available from the last author.

\begin{thebibliography}{00}
\bibitem{disc_o}
K.\ Patel, A.\ Sharma, H.\ Kumar, P.\ Kurundodi, H.\ Burgsteiner, and W.\ Slany,
``A Novel Framework for Decoding Inter-Species Communication with Octopodes
(DISC-O),'' in \textit{OCEANS 2025}, pp.\ 1--9, Sep.\ 2025,
doi: 10.23919/OCEANS59106.2025.11244956.
\bibitem{clip}
A.\ Radford \textit{et al.}, ``Learning Transferable Visual Models from Natural
Language Supervision,'' in \textit{ICML}, 2021.
\bibitem{groundingdino}
S.\ Liu \textit{et al.}, ``Grounding DINO: Marrying DINO with Grounded Pre-Training
for Open-Set Object Detection,'' \textit{arXiv:2303.05499}, 2023.
\bibitem{sam2}
N.\ Ravi \textit{et al.}, ``SAM 2: Segment Anything in Images and Videos,''
\textit{arXiv:2408.00714}, 2024.
\bibitem{lora}
E.\ J.\ Hu \textit{et al.}, ``LoRA: Low-Rank Adaptation of Large Language Models,''
in \textit{ICLR}, 2022.
\bibitem{mobilenetv3}
A.\ Howard \textit{et al.}, ``Searching for MobileNetV3,'' in \textit{ICCV}, 2019.
% TODO: add a Qwen-VL citation and (optionally) an octopus-cognition/welfare reference.
\end{thebibliography}

\end{document}
